Raw multi-axis game scores are dominated by a general quality factor that masks design
distinctiveness: high-scoring games cluster together regardless of how different their
design character actually is. This paper introduces within-game z-score normalization as
the correct method for revealing a game's \emph{design fingerprint} --- the pattern of
what it does above and below its own average across all axes. We show that this
normalization reduces PC1 variance explained from 81.3\% to 50.0\%, surpassing the
pre-registered criterion of $\leq 55\%$, and produces mean pairwise Euclidean distances
of 3.9 (range 1.1--8.7) between game fingerprints, confirming that games genuinely
differ in design character even when they occupy similar quality tiers. We present
fingerprints for 20 illustrative games, demonstrating that the method recovers known
design distinctions (Knizia's elegance emphasis, Lacerda's complexity stack) from
raw scoring data without prior knowledge of designer intent. Codenames emerges as a
singleton cluster with no near neighbours at distance $< 2.5$, confirming its
architectural uniqueness; Pandemic and Chess share a fingerprint cluster despite genre
distance, suggesting structural equivalence masked by thematic framing. Within-game
z-score normalization is the computational foundation for the TIGER framework's claim
to measure distinctiveness rather than quality.
